\documentclass[10pt,a4paper]{article}
\usepackage[margin=25mm]{geometry}
\usepackage{amsmath,amssymb,amsthm,booktabs,mathtools}
\usepackage[T1]{fontenc}
\usepackage{lmodern,microtype}
\setlength{\parskip}{4pt}
\setlength{\emergencystretch}{2em}
\renewcommand{\arraystretch}{1.12}
\usepackage[hidelinks]{hyperref}
\usepackage{xurl}
\newtheorem{theorem}{Theorem}
\newtheorem{definition}[theorem]{Definition}
\newtheorem{lemma}[theorem]{Lemma}
\newtheorem{proposition}[theorem]{Proposition}
\newcommand{\C}{\mathbb C}
\newcommand{\R}{\mathbb R}
\newcommand{\dd}{\mathrm d}
% Generated by build_tables.py from exact bundled certificates.
\newcommand{\MainJetRows}{
$2$ & $-3$ & $4$ & $10$ & $16$ \\
$3$ & $-4$ & $13$ & $-\frac{97}{2}$ & $\frac{1467}{8}$ \\
$\frac{5}{2}$ & $-\frac{7}{2}$ & $\frac{65}{8}$ & $-\frac{391}{32}$ & $\frac{9495}{128}$ \\
}
\newcommand{\MainBoundRows}{
$2$ & $(0.4933991454315, 0.4933991454318)$ & $(-19.70, -19.67)$ & $(-1.01, -0.98)$ \\
$3$ & $(0.4654318384721, 0.4654318384724)$ & $(-432.72, -432.69)$ & $(-13.21, -13.18)$ \\
$\frac{5}{2}$ & $(0.4791575947686, 0.4791575947689)$ & $(-129.86, -129.83)$ & $(-5.17, -5.14)$ \\
}
\newcommand{\MainCoordinateRows}{
$2$ & $(0.972, 0.975)$ & $(0.846, 0.849)$ & $(1.567, 1.570)$ & $(0.848, 0.851)$ & $(7.430, 7.433)$ \\
$3$ & $(0.869, 0.872)$ & $(0.531, 0.534)$ & $(0.820, 0.823)$ & $(0.535, 0.538)$ & $(8.177, 8.180)$ \\
$\frac{5}{2}$ & $(0.918, 0.921)$ & $(0.653, 0.656)$ & $(1.085, 1.088)$ & $(0.655, 0.658)$ & $(7.912, 7.915)$ \\
}
\newcommand{\CertificatePolynomials}{
For $k=2$,
\[\begin{aligned}
f_k(X)&=1152X^{4}-14112X^{3}+59952X^{2}-119970X+46225,\\
g_k(X)&=\frac{12672X^{3}-31392X^{2}-857568X+1565845}{1340525}.
\end{aligned}\]

For $k=3$,
\[\begin{aligned}
f_k(X)&=44000X^{4}-352300X^{3}+1015240X^{2}-1326500X+430923,\\
g_k(X)&=\frac{68200000X^{3}-387643000X^{2}-10844350X+544397116}{867447999}.
\end{aligned}\]

For $k=\frac{5}{2}$,
\[\begin{aligned}
f_k(X)&=38556X^{4}-364248X^{3}+1223067X^{2}-1885095X+660490,\\
g_k(X)&=\frac{3144665916X^{3}-17421405048X^{2}-29056911273X+71805073975}{82898099900}.
\end{aligned}\]
}

\title{Minimal jet orders in a positive behavioral SIR model}
\author{Tatsuhiko Takimoto\\\small Independent researcher}
\date{Version 1.0, 10 November 2026}
\begin{document}
\maketitle
\begin{abstract}
For the behavioral SIR model with transmission rate
$\beta SI/(1+\alpha I)^k$, we determine the least order of an exact
output jet at one time that identifies the rates and hidden initial
states, on the positive real domain with known total population.
For each fixed, known $k\in\{2,3,5/2\}$, generic local identification
requires order four, whereas generic identification against all positive
competitors requires order five, in the sense of identification on an
open dense set of true points. Fourth-order ambiguity is established
on a nonempty open set, without a density claim. The fifth-order upper bound is an
application of the observation-field theorem of Demin, Pogudin and
Rackauckas and a prior identifiability result for this model. To prove
that four derivatives do not suffice globally, we give exact certificates
for two distinct positive points with the same fourth-order jet and
nonzero jet Jacobians at both points. The inverse function theorem
extends this ambiguity to a nonempty open set. The certificates have
rational isolating intervals and can be checked using Python's standard
library. The result concerns exact derivatives, not estimation from
noisy samples.
\end{abstract}

\section{Model and the least-order question}
\label{sec:model}
Consider
\begin{equation}
S'=-\frac{\beta SI}{(1+\alpha I)^k},\qquad
I'=\frac{\beta SI}{(1+\alpha I)^k}-\gamma I,\qquad
R'=\gamma I,\qquad y=I.
\label{eq:model}
\end{equation}
Fix known $k$ and $N_0>0$. The unknown initial point lies in
\[
\Omega_{N_0}=\{(\alpha,\beta,\gamma,S_0,I_0)\in\R_{>0}^5:
S_0+I_0<N_0\},\qquad R_0=N_0-S_0-I_0.
\]
All powers use the positive real branch. Solutions are analytic near
zero. Since the system is autonomous, the observation time can be
translated to zero. Write
\begin{equation}
J_m(\theta)=(I(0),I'(0),\ldots,I^{(m)}(0)).
\label{eq:jet}
\end{equation}
Order $m$ means $m+1$ observed numbers, including the zeroth derivative.

\begin{definition}
\label{def:minimal}
The jet $J_m$ identifies a point $\theta$ locally if there is a
neighborhood $V$ of $\theta$ in $\Omega_{N_0}$ such that
$J_m^{-1}(J_m(\theta))\cap V=\{\theta\}$. It identifies $\theta$
globally if this equality holds with $V=\Omega_{N_0}$.
Generic local or global identifiability means that the corresponding
property holds at every point of an open dense subset of $\Omega_{N_0}$
in the Euclidean topology. In the global definition all competing points
have the same known $k,N_0$, and are not restricted to the generic subset.
\end{definition}

\begin{theorem}
\label{thm:main}
For each fixed, known $k\in\{2,3,5/2\}$ and each known $N_0>0$,
the least jet order for generic local identifiability of
$\alpha,\beta,\gamma,S_0,I_0,R_0$ in~\eqref{eq:model} is four.
The least jet order for generic global identifiability against all
positive competitors is five: $J_4$ contains five observed numbers,
whereas $J_5$ contains six. Moreover, for each of these exponents
there is a nonempty open subset of $\Omega_{N_0}$ whose points have
distinct positive competitors with the same $J_4$.
\end{theorem}
Here $R_0=N_0-S_0-I_0$ is a dependent quantity, so the unknown
parameter-state domain has dimension five.

Pant, Saucedo and Pogudin studied this behavioral model and its
identifiability in~\cite[Section 3.6 and Supplement S9]{PSP,PSPsupp}.
Their supplement provides the auxiliary state
$Z=(1+\alpha I)^k/\beta$. A companion proof~\cite{Companion}
establishes generic identifiability from the complete analytic output
for every fixed known real $k\geq2$. This does not specify the least
finite jet order.

The observation-field results of Demin, Pogudin and Rackauckas
\cite[Theorem 4.12 and Remark 4.18]{DPR} supply the upper bound used
here. Finite-jet reduction and local rank methods are established tools;
related identifiability theory appears in~\cite{HOPY,OPT}. The contribution
of this note is the positive fourth-order ambiguity at two regular
points for three exponents, which makes the known upper bound sharp.
An input--output equation of order two still contains unknown
coefficients, so its order is not the number of derivatives needed to
recover those coefficients and the initial state at one time.

\paragraph{Assurance.}
The proof combines cited mathematical results, the inverse function
and constant rank theorems, and finite exact certificate checks.
Separately written checks support the certificate identities and their
positivity. This is a prose and computer-assisted proof; it has not been
formalized in a proof assistant or reviewed by human experts. The
verification scope and software comparisons are recorded in
\texttt{VERIFICATION-RECORD.md}.

\section{The fifth-order upper bound}
\label{sec:upper}
Write $k=p/q$ in lowest terms. The positive coordinate
$z=(1+\alpha I)^{1/q}$ gives the rational two-state model
\begin{equation}
S'=-\frac{\beta S(z^q-1)}{\alpha z^p},\qquad
z'=\frac{z^q-1}{qz^{q-1}}
       \left(\frac{\beta S}{z^p}-\gamma\right),\qquad
y=\frac{z^q-1}{\alpha}.
\label{eq:cover}
\end{equation}
For $q=1$ ($k=2,3$), this is the original model in the coordinate
$z=1+\alpha I$; no root is taken.
It has two states and three unknown constant parameters. Over $\C$,
let $\mathcal O$ denote its observation field, generated by the output
and all its Lie derivatives. Apply~\cite[Theorem 4.12 and Remark 4.18]{DPR}
with ground subfield $\mathcal E=\C$ and no external inputs. The
states-plus-parameters bound is derivative order five, or the six
observed numbers in $J_5$, hence
\begin{equation}
\mathcal O=\C(y,Dy,\ldots,D^5y).
\label{eq:field}
\end{equation}
Unknown parameter functions are not supplied to $\mathcal E$ as if
they were already observed.

For completeness, the bound in this no-input setting also follows by
a classical differential-field argument. Put
$L=\C(\alpha,\beta,\gamma,S,z)$ and
$F_j=\C(y,Dy,\ldots,D^jy)$. The output $y$ is transcendental over
$\C$, and $D(F_j)\subset F_{j+1}$. Let $r\geq1$ be the first index
for which $D^ry$ is algebraic over $F_{r-1}$. Such an $r\leq5$
exists because $\operatorname{trdeg}(L/\C)=5$. If
$G(T)\in F_{r-1}[T]$ is its minimal polynomial, differentiating
$G(D^ry)=0$ expresses $D^{r+1}y$ in $F_r$: the derivatives of the
coefficients lie in $F_r$, and the separant
$G'(D^ry)\ne0$ in characteristic zero. Hence $D(F_r)\subset F_r$
and all later derivatives lie in $F_r$. Thus
$\mathcal O=F_r=F_5$, proving~\eqref{eq:field} without an additional
finite-generation hypothesis. This explains the cited bound in the
present setting; it is not a new upper-bound result.

The following is the precise cover statement needed here.
\begin{lemma}[Observable quantities on the cover]
\label{lem:observable}
For $(p,q)=(2,1),(3,1),(5,2)$, equip
$L=K(S,z)$, $K=\C(\alpha,\beta,\gamma)$, with the derivation
in~\eqref{eq:cover}, where the five coordinates are algebraically
independent and $D|_K=0$. Its rational constant field is exactly
$\ker D=K$, and $\alpha,\gamma,\beta^q,S,y$ belong to its
observation field $\mathcal O$ over $\C$.
\end{lemma}
The constant-field assertion is the absence of rational first
integrals beyond the parameters, required by
\cite[Theorem 4.7]{OPT}. In the fixed companion version,
Lemma 3 (Section 2) checks that hypothesis; Lemma 4 and Proposition 5
(Section 3) give the minimal normalized input--output equation;
Proposition 6 (Section 4), equations (11)--(13), recovers
$\alpha,\gamma,\beta^q$ from its coefficients. Section 5,
equation (16), recovers $S$. Theorem 2 of that work is the resulting
positive-domain statement; that statement alone is not being used
as a substitute for membership in $\mathcal O$.
Appendix~\ref{app:observable} proves the lemma here, including direct
Wronskian recovery of the coefficients in $\mathcal O$. Thus the
present proof does not need to infer observability from
input--output identifiability via~\cite[Corollary 5.7]{OPT}.
Consequently
\begin{equation}
\alpha,\gamma,\beta^q,S,y\ \in\mathcal O.
\label{eq:observable}
\end{equation}
In particular, for $q=1$ the coefficient formulas recover $\beta$;
for $q=2$ they recover $\beta^2$, not a choice between its complex
square roots. This distinction is retained when using~\eqref{eq:field}.
Positivity is essential to the subsequent global conclusion; no such
conclusion for arbitrary complex or signed parameters is asserted.

\begin{proposition}
\label{prop:upper}
For the three exponents of Theorem~\ref{thm:main}, $J_5$ identifies
a generic true point against every positive competitor at the same
$k,N_0$.
\end{proposition}
\begin{proof}
For each of the five quantities $h$ in~\eqref{eq:observable},
choose a rational representation
$h=A_h(J_5)/B_h(J_5)$ provided by~\eqref{eq:field}.
The pullback of each denominator is nonzero as a rational function.
Let $B$ be their product, and restrict the true point to $B(J_5)\ne0$.
After clearing model denominators, the equalities
$hB_h(J_5)=A_h(J_5)$ are rational identities on the cover.

A positive competitor with the same $J_5$ has exactly the same
nonzero denominator values. The model denominators are also nonzero
there: $\alpha>0$ and $z>1$. Thus the same identities recover
$\alpha,\gamma,\beta^q,S_0,I_0$ at the competitor. Positivity gives a
unique $q$th root $\beta$, and $R_0=N_0-S_0-I_0$ then agrees.

The positive cover with
$S_0+(z_0^q-1)/\alpha<N_0$ is a nonempty real open subset of the
five-dimensional parameter-state space. A nonzero polynomial cannot
vanish on any such open subset. Therefore the nonvanishing condition
is open dense for every $N_0>0$. The positive root substitution is an
analytic diffeomorphism, so its inverse image is open dense in
$\Omega_{N_0}$. No genericity restriction was imposed on the competitor.
\end{proof}
This argument proves the existence of finite rational recovery
expressions on the cover. It does not require their full expansion.
The upper bound is an application of the cited theory, rather than a
new finite-generation theorem.

\section{Exact positive pairs with the same fourth-order jet}
\label{sec:certificates}
It is convenient to use $h=1/Z=\beta/(1+\alpha I)^k$.
The coordinates
\[
(\alpha,\beta,\gamma,S_0,I_0)\longleftrightarrow
(\alpha,\gamma,S_0,y_0,h_0)
\]
are analytically invertible on the positive domain, with inverse
$\beta=h_0(1+\alpha y_0)^k$. The rational vector field becomes
\begin{equation}
S'=-Syh,\qquad y'=Syh-\gamma y,\qquad
h'=-\frac{k\alpha h y'}{1+\alpha y},\qquad
\alpha'=\gamma'=0.
\label{eq:hmodel}
\end{equation}
Thus every finite jet and its Jacobian is rational in these coordinates
for fixed $k$. All Jacobians below use the coordinate order
$(\alpha,\gamma,S_0,y_0,h_0)$.

Put $y=y_0$, $v=I'(0)$, $w=I''(0)$, and define
\begin{equation}
\rho=\frac{\alpha}{1+\alpha y},\qquad
r=v+\gamma y,\qquad F=yw-v^2+k\rho vyr.
\label{eq:F}
\end{equation}
Differentiating $r=Syh$ in~\eqref{eq:hmodel} gives
$F=-y^2rh$. Consequently a positive candidate can be reconstructed as
\begin{equation}
\alpha=\frac{\rho}{1-\rho y},\qquad
h=-\frac{F}{y^2r},\qquad S=-\frac{yr^2}{F},\qquad
\beta=h(1+\alpha y)^k.
\label{eq:reconstruct}
\end{equation}
The conditions are
$0<\rho<1/y$, $\gamma>0$, $r>0$, $F<0$, and $S+y<N_0$.
Here $\rho$ is a continuous coordinate, unrelated to the integer
$q$ in~\eqref{eq:cover}.

For $N_0=10$, take the true point
\begin{equation}
(\alpha,\gamma,S_0,y_0,h_0)=(1,1,2,1,1),
\qquad \beta=2^k,\quad R_0=7.
\label{eq:true}
\end{equation}
In every case $I(0)=I'(0)=1$. The remaining fourth-order data and the
true-point determinant are:
\begin{center}
\begin{tabular}{@{}ccccc@{}}
\toprule
$k$ & $I''(0)$ & $I'''(0)$ & $I^{(4)}(0)$ & $\det DJ_4$ \\
\midrule
\MainJetRows
\bottomrule
\end{tabular}
\end{center}

\begin{lemma}[Exact certificate]
\label{lem:certificate}
For each exponent in Theorem~\ref{thm:main}, there is a distinct
positive point at $N_0=10$ with the same $J_4$ as~\eqref{eq:true}.
Both $5\times5$ Jacobians $DJ_4$ are nonsingular. Their fifth output
derivatives differ.
\end{lemma}
\begin{proof}[Certificate and verification]
Let $f_k$ and $g_k$ be the rational polynomials in
Appendix~\ref{app:data}. Define $\rho_*$ to be the unique real root
of $f_k$ in the displayed $\rho_*$ interval below. Set
$\gamma_*=g_k(\rho_*)$, use $y=v=1$ and the displayed $w$ in
\eqref{eq:F}, and apply~\eqref{eq:reconstruct} to define the second
point. Its removed state is $R_*=9-S_*$. The following are
outward enclosures of the stored exact rational bounds; each decimal
endpoint is itself a rational number.
\begin{center}
\small
\begin{tabular}{@{}cccc@{}}
\toprule
$k$ & $\rho_*$ & $\det DJ_4(\theta_*)$ & $I_*^{(5)}(0)-I^{(5)}(0)$ \\
\midrule
\MainBoundRows
\bottomrule
\end{tabular}
\end{center}

The certificate checker performs these finite exact operations.
Sturm's algorithm counts one root in each displayed interval
(and in its narrower stored interval) and exactly
one root of each $f_k$ in $(0,1)$. The isolating interval excludes
$1/2$. Arithmetic modulo $f_k$
constructs~\eqref{eq:reconstruct}; every inverse used is checked by the
Euclidean algorithm. Rational interval evaluation gives positive lower
bounds for $\alpha_*,\gamma_*,S_*,Z_*=1/h_*,R_*$. In particular the
second point is admissible and distinct from~\eqref{eq:true}.
The coordinate enclosures make this second branch explicit:
\begin{center}
\small
\begin{tabular}{@{}cccccc@{}}
\toprule
$k$ & $\alpha_*$ & $\gamma_*$ & $S_*$ & $Z_*$ & $R_*$ \\
\midrule
\MainCoordinateRows
\bottomrule
\end{tabular}
\end{center}
Here $I_*=1$ in all three rows, while the true coordinates
$(\alpha,\gamma,S,Z,R)$ are $(1,1,2,1,7)$.

Taylor integration of~\eqref{eq:hmodel} through order five, using
ordinary rational coefficients and factorial conversion, gives zero
remainders modulo $f_k$ for all five components of
$J_4(\theta_*)-J_4(\theta)$. Differentiating the same recurrence with
respect to the five initial coordinates gives both Jacobian matrices.
The true determinants are the displayed nonzero rationals. The
alternate determinants and the fifth-derivative differences are
coprime to $f_k$; their interval enclosures also exclude zero.
These computations establish the stated identities and inequalities
at the selected real roots. They do not require trusting a numerical
root finder or a computed Gr\"obner basis.

The data are in \texttt{producer/output/certificate.json}.
Two separately written standard-library checkers use different ODE
coordinates and determinant calculations and agree on the polynomial
remainders. The exact reduced polynomials and interval bounds are
included with the article, so the finite calculation can be repeated.
From the packet directory run \texttt{python3 -B verify.py}; it
replays \texttt{independent/check\_source\_certificate.py} and
\texttt{producer/verify\_certificate.py} with exact arithmetic.
\end{proof}

One route to the candidates is to differentiate
$Z=-y^2r/F$, noting $D\rho=-\rho^2v$. This gives
\begin{equation}
H=yr\,DF-F\{2vr+y(w+\gamma v)-k\rho vyr\}=0.
\label{eq:H}
\end{equation}
For a fourth-order candidate, $H=DH=0$. The coefficient of $I'''$
in $H$, and of $I^{(4)}$ in $DH$, is $y^2r>0$.
For the data in~\eqref{eq:true}, elimination of these two equations
produces a polynomial in $\rho$ containing the known true-point factor
$2\rho-1$. Removing that factor leaves the quartic $f_k$, which
therefore does not include the true point $\rho=1/2$.
The certificate proof above checks the remaining candidate directly
in the original ODE.
It does not use a complete count of all positive candidates.

\section{From regular pairs to least orders}
\label{sec:minimal}
\begin{lemma}
\label{lem:two-branches}
Let $\Phi:U\to\R^d$ be a continuously differentiable map on an open
set $U\subset\R^d$. Suppose distinct points $a,b\in U$ have the same
image and both $D\Phi(a)$ and $D\Phi(b)$ are nonsingular. Then a
nonempty open subset of $U$ consists of points with a distinct
competitor under $\Phi$.
\end{lemma}
\begin{proof}
Take disjoint neighborhoods of $a$ and $b$ on which $\Phi$ is a
diffeomorphism onto open sets. Those two images contain the common
value $\Phi(a)=\Phi(b)$, so their intersection is nonempty and open.
Its inverse image in the first neighborhood is nonempty and open,
and every point there has another inverse image in the second.
\end{proof}

\begin{proof}[Proof of Theorem~\ref{thm:main}]
First take $N_0=10$. The nonzero true determinant in
Lemma~\ref{lem:certificate} shows that the rational function
$\det DJ_4$ is not identically zero. Its nonvanishing set is open
dense in the positive coordinate domain. The inverse function theorem
gives local identifiability there.

For $m\leq3$, $J_m$ maps a five-dimensional open set to dimension
$m+1\leq4$. The maximum-rank locus is open dense: the jet map is
rational in the coordinates of~\eqref{eq:hmodel}, so a nonzero
maximal minor has a zero set with empty interior. The constant rank
theorem gives a positive-dimensional local fiber at every point of
that locus. Such a fiber contains other positive admissible points.
Thus no open dense set can be locally identifiable for $m\leq3$,
and the least generic local order is four.

Apply Lemma~\ref{lem:two-branches} to $\Phi=J_4$ and the two points
of Lemma~\ref{lem:certificate}. The resulting open set of ambiguous
true points prevents any open dense set from satisfying global
identifiability at order four. Lower orders cannot distinguish points
that already share $J_4$. Proposition~\ref{prop:upper} supplies the
order-five upper bound.

Finally let $c=N_0/10>0$. The scaling
\begin{equation}
(S,I,R,Z)\mapsto(cS,cI,cR,cZ),\qquad
(\alpha,\beta,\gamma)\mapsto(\alpha/c,\beta/c,\gamma)
\label{eq:scale}
\end{equation}
preserves~\eqref{eq:model}, multiplies every output derivative by $c$,
and is an analytic diffeomorphism between the two positive parameter
domains. It takes the pair above into $\Omega_{N_0}$, preserving
jet equality and nonsingularity at both points. The same local and
global lower-bound arguments therefore apply for every $N_0>0$.
\end{proof}
The conclusion at order four is an open set of ambiguity. It is not
a claim that every positive point is ambiguous or that every fiber
has exactly two points.

\section{Additional certificates and scope}
\label{sec:scope}
A separate interval calculation uses the polynomial coordinates
\[
\rho=\frac{\alpha}{1+\alpha y},\quad b=\frac{\beta}{(1+\alpha y)^k},
\quad s=Sb,\quad g=\gamma,
\]
with equations
\begin{equation}
\begin{aligned}
y'&=y(s-g),&s'&=-syb-k\rho s y',\\
b'&=-k\rho b y',&\rho'&=-\rho^2y',&g'&=0.
\end{aligned}
\label{eq:polynomial}
\end{equation}
After fixing $y,v,w$, one recovers
$g=s-v/y$ and
$b=-(w/y-v^2/y^2)/(sy)-k\rho v/y$.
Thus the remaining two unknowns $(\rho,s)$ map to
$(I''',I^{(4)})$. On the positive domain the change between
$(\alpha,\gamma,S,y,h)$ and $(y,v,w,\rho,s)$ is invertible,
so nonsingularity of this reduced $2\times2$ Jacobian is equivalent
to nonsingularity of the full $DJ_4$.

The separate supplementary records give interval certificates for
$k=7/2$ and $k=1+\sqrt2$, including the changed true point
$(\alpha,\gamma,S_0,I_0,Z_0,N_0)=(2/3,1,2,1,1,10)$.
The alternate $(\rho,s)$ centers are approximately
$(0.3707239149,1.5505151892)$ and
$(0.3930391625,1.8392705345)$, respectively. The certified
fifth-derivative differences lie in $(-9.752,-9.750)$ and
$(-0.954,-0.952)$. These decimal intervals are outward rational
enclosures, not numerical evidence by themselves.

Let $\mathcal F(\rho,s)$ be the two residuals between the candidate's
third and fourth derivatives and those of the true point. For a rational
box $B$ and its center $a$, the checker forms an invertible rational
preconditioner $C$ and bounds
\[
K(B)=a-C\mathcal F(a)+(\mathrm{Id}-C\mathcal F^{\prime}(B))(B-a).
\]
It verifies $K(B)\subset\operatorname{int}B$ and
$\|\mathrm{Id}-C\mathcal F^{\prime}(B)\|_\infty<1$. The contraction mapping theorem supplies
a unique zero, while interval bounds give positivity, nonsingularity
at both branches, and a nonzero fifth-derivative difference.
For $1+\sqrt2$ the calculation is uniform over a rational interval
whose squared shifted endpoints straddle $2$. The irrational
exponent is not replaced by one rational approximation.
The exact boxes, exponent intervals and rational preconditioners are
stored under each exponent key in
\texttt{independent/stage1-certified-pairs.json} and
\texttt{independent/extra-pairs.json}; the latter contains the changed
initial point. The fields \texttt{root\_box\_p\_q},
\texttt{k\_interval} and \texttt{preconditioner} specify these inputs
(the historical $p,q$ there mean $\rho,s$, not the cover integers).
Run \texttt{python3 -B verify\_intervals.py} to reproduce the seven
stored inclusions and bounds. These supplementary assertions depend
on those separate interval records; the three printed algebraic
certificates suffice for Theorem~\ref{thm:main}.

These additional records corroborate the fourth-order ambiguity
mechanism; the main theorem remains restricted to $2,3,5/2$.
The packet does not classify all exponents or all exceptional points,
and it does not supply an expanded global inverse for $J_5$.
Unknown $k$, unknown $N_0$, numerical stability and estimating
derivatives from noisy observations are outside the result.
Nonsingular jet Jacobians give local inverses; they do not bound
their condition numbers or guarantee stable estimation.

\paragraph{AI-use statement.}
AI systems assisted with the proof development, separate checks,
exact computations, literature comparison and manuscript drafting.
The author is responsible for the mathematical claims and their
stated scope; AI systems are not authors.

\appendix
\section{Polynomials defining the three main certificates}
\label{app:data}
The alternate point is defined by the unique real root $\rho_*$ of
$f_k$ in the displayed interval of Section~\ref{sec:certificates},
and by $\gamma_*=g_k(\rho_*)$. The narrower JSON intervals are
additional implementation data, not part of this definition.
The expressions below are obtained by reducing the recorded
$\gamma$ polynomial modulo $f_k$, without changing its value at a root.
\CertificatePolynomials

\section{A proof of cover observability for the three exponents}
\label{app:observable}
We prove Lemma~\ref{lem:observable}. This is the rational argument of
\cite[Sections 2--5]{Companion}, with a direct coefficient-recovery
step to make the observation-field assertion explicit.

\paragraph{Rational constants.}
The derivation $E=(z^p/y)D$ has the same constants and satisfies
\[
E=-\beta S\partial_S+
 \left(\frac{\alpha\beta}{q}Sz^{1-q}
       -\frac{\alpha\gamma}{q}z^{p-q+1}\right)\partial_z.
\]
For a nonzero constant $f\in K(z)(S)$, expand
$f=\sum_{n\geq m}f_n(z)S^n$ in $K(z)((S))$, with $f_m\ne0$.
The derivation preserves $S$-order, so its lowest coefficient gives
$-\beta m f_m-(\alpha\gamma/q)z^{p-q+1}f'_m=0$.
If $m\ne0$, then $f'_m/f_m$ would have a pole of order
$p-q+1\geq2$ at zero. A rational logarithmic derivative has at
most a simple pole there. Hence $m=0$ and $f_0\in K$.
If $f-f_0$ were nonzero, the same argument would contradict its
positive $S$-order. Thus $\ker D=K$. This verifies the rational
first-integral hypothesis without a numerical test.

\paragraph{The minimal output relation.}
Put $u=1+\alpha y$, $v=Dy$, $w=D^2y$ and
\[
A=yu w-\{1-(k-1)\alpha y\}v^2+k\alpha\gamma y^2v.
\]
Eliminating $S$ from $v+\gamma y=\beta Sy/z^p$ and
$DS=-(v+\gamma y)$, by differentiating the former identity, gives
\begin{equation}
z^{p-q}A=-\beta y^2(v+\gamma y),\qquad
Q=u^{p-q}A^q-(-\beta)^q y^{2q}(v+\gamma y)^q=0.
\label{eq:appQ}
\end{equation}
Here $K(S,z)=K(v,z)$, with $v,z$ independent; thus $y,v$ are
independent. Over $F=K(y,v)$, $T^q-u$ is Eisenstein at $u$
in $K(v)[u][T]$ (or linear for $q=1$).
The expression for $w$ in~\eqref{eq:appQ} is an element of $F$
plus a nonzero element of $F$ times $z^{-(p-q)}$.
Since $\gcd(p-q,q)=1$, its $q$ conjugates are distinct.
Therefore $Q$ has minimal degree $q$ in $w$ and is irreducible
over $F$. Its leading coefficient in $w$ is $y^q u^p$;
the only possible content factors in $K[y,v][w]$ are $y,u$.
But $Q(0,v,w)=(-1)^qv^{2q}$ and
$Q(-1/\alpha,v,w)=-(-\beta)^q\alpha^{-2q}(v-\gamma/\alpha)^q$
are nonzero. Hence $Q$ is primitive and irreducible in $K[y,v,w]$.
The kernel of evaluation $K[y,v,w]\to L$ is exactly $(Q)$:
division over $F$ and primitivity give this assertion.
The monomial coefficient $[y^qw^q]Q$ is one.

\paragraph{Coefficients are in the observation field.}
Let $m_0=y^qw^q,m_1,\ldots,m_n$ be the monomial support of $Q$
over $K$, evaluated in $L$, and write $Q=m_0+\sum_{j=1}^n c_jm_j$.
A $K$-linear relation among these evaluations lies in $(Q)$.
Its total degree is at most $\deg Q$, so it is a scalar multiple
of $Q$. In particular $m_1,\ldots,m_n$ are independent over $K$.
Since $\ker D=K$, the characteristic-zero Wronskian criterion gives
$W=\det(D^im_j)_{0\leq i<n,\,1\leq j\leq n}\ne0$ in $L$.
For clarity, the criterion follows by induction: divide by a nonzero
first element and use
$W(f_1,\ldots,f_n)=f_1^nW(1,f_2/f_1,\ldots,f_n/f_1)$;
expansion along the first column reduces to the derivatives of the
ratios. A constant relation among those derivatives integrates to
a constant relation among the original elements. The converse is
obtained by differentiating a constant relation.

Differentiate $Q=0$ for $i=0,\ldots,n-1$. Since $Dc_j=0$,
the invertible system
$\sum_j(D^im_j)c_j=-D^im_0$ expresses each $c_j$ as a rational
function of output derivatives. Thus every coefficient of $Q$
belongs to $\mathcal O$. This is a field identity, not merely
equality of identifiable values on positive real solutions.
Although this construction uses derivatives of possibly high order,
\eqref{eq:field} subsequently places the same quantities in $F_5$;
there is no assumed fifth-order inverse in this argument.

\paragraph{Parameter and state recovery.}
For $q=1$ and $p=2,3$, expansion of~\eqref{eq:appQ} gives
\[
[y^2w]Q=p\alpha,\quad
[y^3v]Q=p(p-1)\alpha^2\gamma,\quad
[y^2v]Q=\beta+p\alpha\gamma.
\]
These recover $\alpha,\gamma,\beta$ successively in $\mathcal O$.
For $(p,q)=(5,2)$, the needed coefficients are
\[
[y^3w^2]Q=5\alpha,\quad
[y^3vw]Q=5\alpha\gamma,\quad
[y^6]Q=-\beta^2\gamma^2.
\]
They recover $\alpha,\gamma,\beta^2$ in $\mathcal O$.
All divisions are by nonzero elements of the rational function field.
Finally~\eqref{eq:appQ} and $v+\gamma y=\beta Sy/z^p$ imply
\begin{equation}
S=-\frac{yu(v+\gamma y)^2}{A}\in\mathcal O.
\label{eq:appS}
\end{equation}
The denominator is not identically zero; along a positive solution
$z^{p-q}A=-\beta^2Sy^3/z^p\ne0$. Also $y\in\mathcal O$ by
definition. This completes the proof of the lemma.

\clearpage
\begin{thebibliography}{9}
\bibitem{PSP}
B. Pant, O. Saucedo and G. Pogudin,
``Uncovering identifiability of epidemiological models: basic reproduction
number and complementary data streams,'' medRxiv, version 2
(25 January 2026), doi:\href{https://doi.org/10.64898/2026.01.16.26344284}{10.64898/2026.01.16.26344284}.
\bibitem{PSPsupp}
B. Pant, O. Saucedo and G. Pogudin,
Supplementary material to~\cite{PSP}, version 2, Section S9,
\url{https://www.medrxiv.org/content/medrxiv/early/2026/01/25/2026.01.16.26344284/DC1/embed/media-1.pdf}.
\bibitem{DPR}
A. Demin, G. Pogudin and C. Rackauckas,
``Observable functions of rational ODE models and how to find them,''
\href{https://arxiv.org/abs/2609.06134v3}{arXiv:2609.06134v3}
(18 September 2026). The numbering used here is that version's.
\bibitem{Companion}
T. Takimoto,
``Global identifiability of a behavioral SIR model for every fixed real
saturation exponent $k\geq2$,'' article version 3.0
(27 October 2026), publication display edition (10 November 2026),
Lemma 3, Lemma 4, Propositions 5--6 and Section 5, equations (14)--(16).
The display edition is included in the \texttt{companion/} directory
of the accompanying packet. The TeX source has SHA-256
\texttt{6689d03fc3aaeeb9741b60c4169c831e5ecb3b20eb1884a36baa09985d940d3b};
the PDF has SHA-256
\texttt{ab6e2378648e9da9972c19ec79ed9a32c9503bcf8b50f27020cf9d63011b4e14}.
\bibitem{OPT}
A. Ovchinnikov, G. Pogudin and P. Thompson,
``Parameter identifiability and input--output equations,''
\emph{Appl. Algebra Engrg. Comm. Comput.} \textbf{34}(2), 165--182 (2023),
doi:10.1007/s00200-021-00486-8; arXiv:2007.14787v2.
The theorem numbering used here is that of the cited arXiv version.
\bibitem{HOPY}
H. Hong, A. Ovchinnikov, G. Pogudin and C. Yap,
``Global identifiability of differential models,''
\emph{Comm. Pure Appl. Math.} \textbf{73}(9), 1831--1879 (2020),
doi:10.1002/cpa.21921; arXiv:1801.08112v5.
\end{thebibliography}
\end{document}
